% ECONOMICS_PROBLEM: {"id": "L41-475", "title": "Serial acquisitions", "jel_code": "L41", "source_id": "OP-0395"}
% !TeX program = xelatex
\documentclass[11pt,a4paper]{article}
\usepackage[margin=23mm]{geometry}
\usepackage{fontspec}
\setmainfont{Latin Modern Roman}
\usepackage{amsmath,amssymb,mathtools}
\usepackage{xcolor,enumitem,booktabs,longtable,array,xurl,hyperref}
\definecolor{muted}{HTML}{53636C}
\hypersetup{colorlinks=true,urlcolor=blue,linkcolor=blue}
\setlength{\parindent}{0pt}
\setlength{\parskip}{5pt}
\newcommand{\R}{\mathbb R}
\newcommand{\E}{\mathbb E}
\newcommand{\N}{\mathbb N}
\newcommand{\ind}{\mathbf 1}
\newcommand{\Var}{\operatorname{Var}}
\newcommand{\Cov}{\operatorname{Cov}}
\newcommand{\supp}{\operatorname{supp}}
\DeclareMathOperator*{\argmin}{arg\,min}
\DeclareMathOperator*{\argmax}{arg\,max}
\newcommand{\entrymeta}[3]{{\sffamily\small\color{muted}\textbf{\texttt{#1}}\quad|\quad #2\par #3\par}}
\newcommand{\Problemref}[1]{\texttt{#1}}
\newcommand{\JELref}[2]{\texttt{#2} (JEL \texttt{#1})}
\begin{document}
\section*{Serial acquisitions}
\textbf{Problem ID:} \texttt{L41-475}\quad \textbf{JEL:} \texttt{L41}\par
% BEGIN ECONOMICS STATEMENT
\label{problem:OP-0395}
{\sffamily\small\color{muted}Original subject group: Industrial organization and competition\par}
\label{definitions:problem:30007677}
\textbf{Audit classification:} Background; proposed extension. The source verifies a policy concern about serial acquisitions; it does not publish this formal cumulative-price question. Isolated-deal effects must be distinguished from telescoping sequential increments.
\subsection*{Definitions and precise research target}
Fix local United States outpatient-provider markets and a ten-year observation window. For each acquirer \(m\), let \(A_m=(a_{m1},\ldots,a_{mK_m})\) be its observed sequence of acquisitions, where \(K_m\) is the number of deals and every deal is individually below the contemporaneous notification threshold. The threshold is a descriptive selection rule, not an assumption of competitive harmlessness. Define \(P_m(A_m,h)\) as the equilibrium price index at horizon \(h=5\) years under that acquisition sequence, permitting rival entry, exit and investment. Define \(P_m(0,h)\) under the counterfactual sequence containing no acquisitions by that acquirer, with other institutional changes specified identically. The target cumulative effect is \(\tau(h)=E[P_m(A_m,h)-P_m(0,h)]\), averaged over this below-threshold acquirer population. Compare it with the sum of effects attributed to the individual deals to measure nonadditivity. Construct ownership histories covering unreported transactions and address endogenous choice of targets and acquisition timing. A valid comparison requires stated counterfactual-trends or instrument assumptions, with bounds when these assumptions are uncertain. Report heterogeneity by baseline concentration and overlapping service lines. This agenda quantifies cumulative effects in one service sector; it does not infer an effect from notification status alone or claim that all small acquisitions reduce competition.

\textbf{Definitions, assumptions and mathematical qualifications.} The input sequence \(A_m\) includes each acquired target, closing date and acquired service line. Its zero counterpart holds all other policy rules fixed but permits endogenous responses. Define isolated-deal sequence \(A_m^{(k)}\) as only deal \(k\), at its observed date, and use a common terminal calendar date. The nonadditivity contrast is \(\mathcal N_m=P_m(A_m,h)-P_m(0,h)-\sum_k\{P_m(A_m^{(k)},h)-P_m(0,h)\}\). This is different from the sum of sequential marginal effects, which telescopes to the total by construction. Eligibility uses the notification rule in force on each actual transaction date; historical thresholds must be independently sourced rather than replaced by today's threshold.
\subsection*{Variants and assumption changes}
\begin{enumerate}
\item \textbf{Editorial, additive benchmark.} Under \(P_m(A)=P_m(0)+\sum_{k\in A}\gamma_{mk}\), \(\mathcal N_m=0\) algebraically; estimate departures rather than call additivity a general open theorem.
\item \textbf{Editorial, endogenous histories.} Model acquisition timing and targets as part of a dynamic structural law or use a specified quasi-experiment. Report bounds over admissible histories when observed transactions and selection mechanisms are incomplete.
\end{enumerate}
\subsection*{References and source audit}
\textbf{Checked source.} 24 August 2024 author draft, Section 6, printed p. 35, serial-acquisition policy recommendation. Full primary text retrieved. \url{https://web.stanford.edu/~ayurukog/trendsincompetition.pdf}.
\begin{enumerate}\raggedright
\item\hypertarget{definitions:source:30007677:1}{} Carl Shapiro; Ali Yurukoglu. 2024. \emph{Trends in Competition in the United States: What Does the Evidence Show?}. 24 August 2024 author draft, Section 6, printed p. 35, serial-acquisition policy recommendation.
\url{https://web.stanford.edu/~ayurukog/trendsincompetition.pdf}.
DOI: \url{https://doi.org/10.3386/w32762}.
\end{enumerate}

\subsection*{Common conventions: Source mathematical conventions}

These conventions apply to each entry unless explicitly replaced.

\textbf{Models and identification.} A primitive model \(m\) specifies all probability laws, feasible choices, information, equilibrium and measurement rules used by its target. The admissible class \(\mathcal M\) is fixed independently of outcomes. If \(P_O(m)\) is its observable probability law and \(\tau(m)\) the target, its identified set at observable law \(P\) is
\[
 \mathcal I_\tau(P)=\{\tau(m):m\in\mathcal M,\ P_O(m)=P\}.
\]
Point identification means that this set is a singleton. An empty set signals incompatibility of data and assumptions. Coordinate bounds need not describe the joint set. Bounds are sharp only if every retained target is attainable or arbitrarily approachable by compatible models. Finite bounds require explicit restrictions on unobserved quantities; unrestricted confounding, hidden wealth, extrapolation or arbitrary equilibrium selection can yield uninformative sets.

\textbf{Causal interventions.} A potential outcome \(Y_i(p)\) is the outcome under a complete policy regime \(p\), including timing, financing, information and market spillovers. The target population and units are fixed before treatment. With interference, use the complete assignment vector or a justified exposure mapping; individual no-interference notation alone cannot identify an economy-wide intervention. A difference of mean potential outcomes does not require the cross-world joint distribution. Conditioning on post-treatment migration, survival, enrollment or employment changes the estimand and needs separate justification.

\textbf{Statistical statements.} An estimator, confidence set, forecast or test specifies a sampling law, model class and loss/coverage criterion. Uniform \((1-\alpha)\)-coverage means \(\inf_{m\in\mathcal M}\Pr_m\{\tau(m)\in C_n\}\geq1-\alpha\), or an explicitly asymptotic version. The whole parameter space trivially has coverage, so informativeness requires a stated width, risk or power objective. A forecasting improvement is relative to a named benchmark, information set and evaluation population.

\textbf{Equilibrium and optimization.} An equilibrium comprises feasible plans, optimality, beliefs where needed, budget constraints and all market/resource clearing. The primitives or a stated benchmark define each correspondence \(\mathcal E(p;m)\). Nonemptiness is assumed or proved, never inferred just from compact policy sets. A selected-equilibrium target uses a declared selection rule; a robust target quantifies over all permitted equilibria. Use suprema or infima unless attainment is established. An \(\varepsilon\)-optimal policy is feasible and within \(\varepsilon\) of the finite optimum value. Report an empty feasible set separately.

\textbf{Welfare and accounting.} Normative utility, interpersonal normalization, social weights, numeraire, discounting and the use of public receipts are primitives. Behavioral utility need not coincide with the welfare criterion. Household welfare includes ownership income and rents once; adding profits again to compensating variation generally double-counts them. A monetary equivalent is defined by an explicit transfer and price convention, with monotonicity and range conditions. Average effects, mechanism contributions, transfers and welfare losses are different targets. Infinite sums require convergence and dynamic feasible plans require terminal or no-Ponzi conditions.

\textbf{Source versions.} A working-paper date, revision date and journal publication year are distinguished. A DOI for a journal version is not automatically the DOI of an earlier manuscript. Locators refer to the stated version and printed pages where specified. A metadata-only check does not verify a claim about a full-text conclusion. Every entry's source subsection narrows the provenance claim accordingly.
% END ECONOMICS STATEMENT
\end{document}
